% P47の固定条件に対応する方法部品。中心的新規性・完成原稿ではない。
\subsection{仕様の選択表示と拒否時の契約返信}
仕様表示を省略した後の修復費用を切り分けるため，
全表示と詳しい契約返信，選択表示と引数名の返信，
選択表示と詳しい契約返信の3条件を比較する．
既知の開発用4課題について第8応答後の同じ会話・環境・Python状態から
各条件を1回実行し，全12条件を集計する．
選択表示は，保存された表示式とAPI返却・出力全体の由来が一致する
古い仕様部分だけを一度省略する．その後の出力は通常どおり追記する．
全表示にも同じ詳しい契約返信を与えることで，
修復情報を片側だけに追加した効果を表示削減の効果へ含めない．
引数名だけの返信を，詳細契約の必要性に対する単純対照とする．

返信の対象は，実行中のプログラムが要求したAPIのうち，
契約が既に公開され，未知引数により送信前に拒否された要求に限る．
詳しい返信は元の公開契約をそのまま用い，必須指定や既定値を変更しない．
返信の文字数が16,000を超える場合は引数名へ，
それも超える場合は元の拒否へ戻す．
自動的な引数修正や再送は行わず，
同じプログラム内で既に成功した作用とPython状態を保持する．
追加の生成と操作はモデルの後続行動として費用へ含める．

各条件の継続上限を16生成，入力と出力の合計120,000トークン，480秒とし，
API上限は共通前半を含む元課題の値とする．
残予算通知，残量による出力枠，初回キャッシュ不使用は全条件に共通とする．
独立評価の成否，既知使用量と不明呼出し，API，時間，制御費用を記録する．
共通前半を加えた論理的な全仕事の費用と，実験用復元・ロード費用を分ける．
使用量や作用が不明の場合などは事前規則で全体を停止し，
未評価・未実行を成功や費用ゼロの完遂へ置き換えない．

原要求と返却を照合し，追加された契約が元の公開範囲に含まれること，
送信済み要求へ適用されていないこと，後続入力との対応を監査する．
全12条件の評価・既知使用量・監査合格を前提に，選択表示と詳しい返信が
両対照の成功を課題ごとに失わず，それぞれの継続総トークンを下回ることを
開発上の採用基準とする．4既知課題・各1回のため，
この基準を一般的な品質維持，統計的非劣性や方式の新規性の証明とはしない．
